\begin{abstract}
Corrupted training data can remove correct evidence and teach false class
relationships. I compared clean training with three forms of 10--30\%
contamination: mixed target-class errors, a fixed cyclic labeling error, and
images from classes absent from the task. Each of the ten data conditions was
trained ten times, giving 100 classifiers. Every classifier was evaluated on
ten balanced test sets. At 30\%, mean accuracy was 73.65\% for mixed noise,
70.03\% for systematic noise, and 90.17\% for unseen-class contamination,
compared with 96.22\% for clean training. Systematic noise caused the largest
loss at 10\% and 30\%, while mixed noise was 0.30 percentage points lower at
20\%. The systematic models also reproduced the imposed error directions.
Thus, contamination rate mattered, but the relationship between true content
and assigned label mattered more than the presence of irrelevant images alone.
\end{abstract}
